\section{Clone a Linked List with Random Pointers}
\label{sec:ll-clone-random}

\problemstatement{Each node has a \texttt{next} pointer and a
\texttt{random} pointer to any node (or null). Produce a deep copy of the
list -- new nodes with correctly wired \texttt{next} and \texttt{random}
pointers -- in $O(1)$ extra space.}

\begin{patternbox}
The clever $O(1)$-space method \textbf{interleaves} each copy right after
its original. Then a copy's \texttt{random} is just the original's random's
\texttt{next}. A final pass unweaves the two lists apart.
\end{patternbox}

\subsection*{Interleave, wire randoms, unweave}

\textbf{Pass 1:} for each original node $x$, create a copy $x'$ and insert it
between $x$ and $x.\texttt{next}$, forming
$x\to x'\to x.\texttt{next}\to\ldots$ \textbf{Pass 2:} set each copy's random
via $x'.\texttt{random}=x.\texttt{random}.\texttt{next}$ -- because every
original is immediately followed by its copy, the copy of any target is one
\texttt{next} away. \textbf{Pass 3:} separate the interleaved list back into
the original and the cloned list, restoring \texttt{next} pointers.

\begin{ideabox}[The Copy Sits Next to Its Original]
Interleaving means ``the clone of node $x$'' is always \texttt{x->next}. So
the random pointer of a clone is \texttt{x->random->next} -- computable
without a hash map, which is what removes the $O(n)$ auxiliary space a map-
based clone would use.
\end{ideabox}

\subsection*{Algorithm}

\begin{enumerate}
  \item Insert each node's copy immediately after it.
  \item For each copy, set \texttt{copy->random = orig->random->next} (guard
        null).
  \item Detach the copies into a separate list, restoring the originals'
        \texttt{next}.
\end{enumerate}

\subsection*{Dry run}

\dryrun{$A\to B$ with $A.\texttt{random}=B$, $B.\texttt{random}=B$.}

\begin{itemize}
  \item Interleave: $A\to A'\to B\to B'$. Set $A'.\texttt{random}=
        A.\texttt{random}.\texttt{next}=B'$; $B'.\texttt{random}=B'$.
  \item Unweave: clone is $A'\to B'$ with correct randoms.
\end{itemize}

\subsection*{C++ implementation}

\begin{lstlisting}
#include <bits/stdc++.h>
using namespace std;

struct Node {
    int val;
    Node* next;
    Node* random;
    Node(int v) : val(v), next(nullptr), random(nullptr) {}
};

Node* copyRandomList(Node* head) {
    if (!head) return nullptr;
    // Pass 1: interleave copies
    for (Node* c = head; c; c = c->next->next) {
        Node* copy = new Node(c->val);
        copy->next = c->next;
        c->next = copy;
    }
    // Pass 2: wire random pointers
    for (Node* c = head; c; c = c->next->next)
        if (c->random) c->next->random = c->random->next;
    // Pass 3: unweave
    Node dummy(0), *t = &dummy;
    for (Node* c = head; c; c = c->next) {
        t->next = c->next;                        // pick the copy
        t = t->next;
        c->next = c->next->next;                  // restore original
    }
    return dummy.next;
}

int main() {
    Node* a = new Node(1);
    Node* b = new Node(2);
    a->next = b;
    a->random = b; b->random = b;
    Node* cl = copyRandomList(a);
    for (Node* c = cl; c; c = c->next)
        cout << c->val << "(rand " << (c->random ? c->random->val : -1) << ") ";
    cout << "\n";                                  // 1(rand 2) 2(rand 2)
    return 0;
}
\end{lstlisting}

\begin{complexitybox}
\textbf{Time Complexity.} $O(n)$ -- three linear passes. \textbf{Space
Complexity.} $O(1)$ extra (versus $O(n)$ for a hash-map clone).
\end{complexitybox}

\begin{takeawaybox}
Cloning with random pointers in $O(1)$ space uses interleaving so each
clone sits beside its original, making random wiring a local
\texttt{->next} lookup. It is the standout example of exploiting the list's
own structure instead of an auxiliary map.
\end{takeawaybox}
